\section{Definite integrals and accumulation}

An indefinite integral gives a family of antiderivatives. A definite integral gives a number. It measures accumulated signed area over an interval.

The notation is
\[
\int_a^b f(x)\,dx.
\]
The numbers \(a\) and \(b\) are the limits of integration. The function \(f\) is the integrand. The variable \(x\) is the variable of integration.

\subsection*{Area under a non-negative function}

If \(f(x)\ge0\) on \([a,b]\), then
\[
\int_a^b f(x)\,dx
\]
can be interpreted as the area under the graph of \(f\), above the \(x\)-axis, from \(x=a\) to \(x=b\).

\begin{examplebox}
Let
\[
f(x)=2
\]
on \([1,4]\). The graph is a horizontal line. The region under the graph is a rectangle of height \(2\) and width \(3\). Therefore
\[
\int_1^4 2\,dx=2\cdot3=6.
\]
\end{examplebox}

This example is simple, but it gives the right idea: a definite integral accumulates contributions over an interval.

\subsection*{Signed area}

If the function is below the \(x\)-axis, then the definite integral contributes negatively. Therefore a definite integral is not always ordinary geometric area. It is signed area.

\begin{examplebox}
Let
\[
f(x)=-1
\]
on \([0,3]\). The graph lies below the \(x\)-axis. The signed area is
\[
\int_0^3 (-1)\,dx=-3.
\]
The ordinary geometric area of the region is \(3\), but the integral is \(-3\).
\end{examplebox}

This distinction matters. If a function is sometimes positive and sometimes negative, the integral can involve cancellation.

\subsection*{Basic properties}

Definite integrals satisfy natural rules:
\[
\int_a^b (f(x)+g(x))\,dx=\int_a^b f(x)\,dx+\int_a^b g(x)\,dx,
\]
\[
\int_a^b c f(x)\,dx=c\int_a^b f(x)\,dx.
\]
They are also additive over intervals:
\[
\int_a^c f(x)\,dx=\int_a^b f(x)\,dx+\int_b^c f(x)\,dx.
\]

Changing the order of the limits changes the sign:
\[
\int_b^a f(x)\,dx=-\int_a^b f(x)\,dx.
\]
Also,
\[
\int_a^a f(x)\,dx=0.
\]
There is no interval over which to accumulate.

\subsection*{Definite integrals are numbers}

The result of a definite integral is a number, not a function of \(x\). The variable of integration is internal to the integral.

For example,
\[
\int_0^2 x^2\,dx
\]
is a number. After integration, the variable \(x\) no longer remains.

\begin{remarkbox}
In an indefinite integral, the answer is a family of functions and includes \(+C\). In a definite integral, the answer is a number and no \(+C\) remains.
\end{remarkbox}

\begin{summarybox}
When reading a definite integral, ask:
\begin{itemize}
\item What interval is being used?
\item Is the function positive, negative, or changing sign?
\item Is the integral measuring signed area or ordinary area?
\item Is the answer expected to be a number or a function?
\item Can the integral be split into simpler intervals?
\end{itemize}
\end{summarybox}

A definite integral measures accumulation across an interval. It is one of the main ways calculus turns local quantities into total quantities.